\documentclass[10pt]{article}

\usepackage[margin=1in]{geometry}
\usepackage{amsmath,amssymb,bm,mathtools}
\usepackage{algorithm}
\usepackage{algpseudocode}
\usepackage{booktabs}
\usepackage{graphicx}
\usepackage{float}
\usepackage{microtype}
\usepackage[hidelinks]{hyperref}

\newcommand{\Mink}{\langle\cdot,\cdot\rangle_M}
\newcommand{\HH}{\mathcal{H}}
\newcommand{\Exp}{\operatorname{Exp}}
\newcommand{\Log}{\operatorname{Log}}
\newcommand{\Proj}{\Pi}
\newcommand{\dist}{d_m}
\newcommand{\R}{\mathbb{R}}

\title{JetFUEL: Jet Flow Matching with Equivariance\
Under the Lorentz Group}
\author{}
\date{}

\begin{document}
\maketitle

\begin{abstract}
We introduce JetFUEL (Jet Flow Matching with Equivariance Under the Lorentz Group),
a Lorentz-equivariant message-passing vector field trained by
Riemannian flow matching on a regulated mass shell. The network uses evolving
auxiliary Lorentz vectors for message passing, while its output field is built
intrinsically from mass-shell logarithmic maps and normalized tangent projections
of two condition-derived reference vectors. This document specifies the geometry,
network, conditional prior, coupling, objective, and sampling procedure.
\end{abstract}

\section{Geometry and notation}

Four-vectors are ordered as $p=(p^0,p^1,p^2,p^3)=(E,p_x,p_y,p_z)\in\R^{1,3}$
and use metric signature $(+,-,-,-)$,
\begin{equation}
 \langle p,q\rangle_M=p^0q^0-\sum_{a=1}^3p^aq^a,
 \qquad \|p\|_M^2\equiv p^2=\langle p,p\rangle_M.
\end{equation}
For a fixed radius (regulator mass) $m>0$, the forward hyperboloid is
\begin{equation}
 \HH_m=\{p\in\R^{1,3}:p^2=m^2,\ p^0>0\}.
\end{equation}
This is the radius-$m$ version of the unit hyperboloid used by
Lou et al.~\cite[Table~2]{lou2020neural}; their convention is obtained with
$x=p/m$ and $\langle x,y\rangle_L=-\langle p,q\rangle_M/m^2$.

The tangent space and induced Riemannian metric of this hyperboloid are
the standard expressions~\cite[Table~2]{lou2020neural}:
\begin{equation}
 T_p\HH_m=\{u:\langle p,u\rangle_M=0\},\qquad
 g_p(u,v)=-\langle u,v\rangle_M,\qquad
 \|u\|_p=\sqrt{-\langle u,u\rangle_M}.
\end{equation}
The corresponding ambient-to-tangent projection is~\cite[Table~2]{lou2020neural}
\begin{equation}
 \Proj_p(v)=v-\frac{\langle p,v\rangle_M}{m^2}p.
 \label{eq:projection}
\end{equation}

For $p,q\in\HH_m$, the hyperbolic distance is~\cite[Table~2]{lou2020neural}
\begin{equation}
 \gamma(p,q)=\frac{\langle p,q\rangle_M}{m^2},\qquad
 \dist(p,q)=m\operatorname{arcosh}\!\bigl(\gamma(p,q)\bigr).
\end{equation}
The associated logarithmic and exponential maps are~\cite[Table~2]{lou2020neural}
\begin{align}
 \Log_p(q)&=
 \frac{\operatorname{arcosh}(\gamma(p,q))}
 {\sinh(\operatorname{arcosh}(\gamma(p,q)))}
 \bigl(q-\gamma(p,q)p\bigr),
 \label{eq:log-map}\\
 \Exp_p(u)&=\cosh\!\left(\frac{\|u\|_p}{m}\right)p+
 m\sinh\!\left(\frac{\|u\|_p}{m}\right)\frac{u}{\|u\|_p}.
 \label{eq:exp-map}
\end{align}
Both expressions use their continuous values at $q=p$ and $u=0$.

\begin{center}
\fbox{\begin{minipage}{0.91\linewidth}
\small\textbf{Geometric interpretation.}
With $w_{p\to q}=q-\gamma(p,q)p$ and
$\|w_{p\to q}\|_p=m\sinh(\operatorname{arcosh}\gamma(p,q))$, the same map is
$\displaystyle \Log_p(q)=\dist(p,q)w_{p\to q}/\|w_{p\to q}\|_p$.
Thus the log map is the unit tangent direction from $p$ toward $q$, scaled by
the geodesic distance.
\end{minipage}}
\end{center}

A jet with $N$ real constituents is modelled as a point on the standard product
manifold.  Its product metric and componentwise maps are likewise standard
\cite{chen2023flow}:
\begin{equation}
 \mathcal{M}_N=\HH_m^N,\qquad
 G_X(U,V)=\sum_{i=1}^N g_{x_i}(u_i,v_i),
 \label{eq:product-metric}
\end{equation}
where $X=(x_1,\ldots,x_N)$ and $U=(u_1,\ldots,u_N)\in T_X\mathcal{M}_N$.
The product maps act componentwise:
\begin{equation}
 \Exp_X(U)=(\Exp_{x_1}(u_1),\ldots,\Exp_{x_N}(u_N)),\qquad
 \Log_X(Y)=(\Log_{x_1}(y_1),\ldots,\Log_{x_N}(y_N)).
 \label{eq:product-maps}
\end{equation}
The choice to use this product manifold for jet constituents, and the masking of
padded slots in implementation, are model-specific conventions.

\section{Conditional inputs}

The vector field receives the current physical state $X\in\mathcal{M}_N$, flow
time $t\in[0,1]$, and condition $c=[c_{\rm type},N,p_{T,J},m_J]$. It also receives
the ordered references $r_0=e_t=(1,0,0,0)$ and $r_1=p_J$. The latter is constructed
from the conditioned jet attributes,
\begin{equation}
 p_J=\frac{1}{s}
 \begin{pmatrix}
 m_{T,J}\cosh\eta_J\\
 p_{T,J}\cos\phi_J\\
 p_{T,J}\sin\phi_J\\
 m_{T,J}\sinh\eta_J
 \end{pmatrix},\qquad m_{T,J}=\sqrt{p_{T,J}^2+m_J^2},
 \label{eq:jet-reference}
\end{equation}
where $s$ is the dataset momentum scale. These vectors are condition-derived;
they are not formed by summing the target constituents.

\section{JetFUEL vector field}

The network maintains invariant scalar features $h_i^\ell\in\R^{d_h}$ and
auxiliary Lorentz-vector features $y_i^\ell\in\R^{1,3}$, with $d_h=96$ and
$L=6$ layers. The physical state $x_i$ remains fixed during a network evaluation.
Only $y_i^\ell$ evolves through message passing. The graph is complete over real
constituents, with $\mathcal{N}(i)=\{j:j\neq i\}$.

\begin{algorithm}[H]
\caption{JetFUEL vector field $V_\theta(X,t,c;r_0,r_1)$}
\label{alg:network}
\small
\begin{algorithmic}[1]
\Require $X=(x_1,\ldots,x_N)$, $t$, $c$, $(r_0,r_1)$
\Ensure $V_\theta(X,t,c;r_0,r_1)\in T_X\mathcal{M}_N$
\State $y_i^0\gets x_i$ and $\tau(t)\gets[t,\sin(\pi t),\cos(\pi t)]$ for all $i$
\State $h_i^0\gets\phi_{\rm seed}([\psi(\langle y_i^0,r_0\rangle_M),\psi(\langle y_i^0,r_1\rangle_M)])+\phi_c(c)-\phi_t(\tau(t))$
\For{$\ell=0,\ldots,L-1$}
  \State $\bar h_i^\ell\gets\operatorname{LayerNorm}(h_i^\ell)$
  \For{each $j\in\mathcal{N}(i)$}
    \State $\xi_{ij}^\ell\gets[\psi(\|y_i^\ell-y_j^\ell\|_M^2),\psi(\langle y_i^\ell,y_j^\ell\rangle_M)]$
    \State $M_{ij}^\ell\gets\phi_e^\ell(\bar h_i^\ell,\bar h_j^\ell,\xi_{ij}^\ell)$; \quad $\alpha_{ij}^\ell\gets\phi_x^\ell(M_{ij}^\ell)$
  \EndFor
  \State $h_i^{\ell+1}\gets h_i^\ell+\phi_h^\ell\!\left(\bar h_i^\ell,|\mathcal{N}(i)|^{-1/2}\sum_{j\in\mathcal{N}(i)}M_{ij}^\ell\right)$
  \State $y_i^{\ell+1}\gets y_i^\ell+\varepsilon_y|\mathcal{N}(i)|^{-1/2}\sum_{j\in\mathcal{N}(i)}\alpha_{ij}^\ell(y_i^\ell-y_j^\ell)$
\EndFor
\For{each $j\in\mathcal{N}(i)$}
  \State $\xi_{ij}^L\gets[\psi(\|y_i^L-y_j^L\|_M^2),\psi(\langle y_i^L,y_j^L\rangle_M)]$
  \State $a_{ij}\gets\phi_f(h_i^L,h_j^L,\xi_{ij}^L)$; \quad $q_{ij}\gets\Log_{x_i}(x_j)/(m+\dist(x_i,x_j))$
\EndFor
\State $v_i^{\rm part}\gets|\mathcal{N}(i)|^{-1/2}\sum_{j\in\mathcal{N}(i)}a_{ij}q_{ij}$
\State $(b_{i0},b_{i1})\gets\phi_r(h_i^L,[\psi(\langle y_i^L,r_0\rangle_M),\psi(\langle y_i^L,r_1\rangle_M)])$
\State $q_{ir}^{\rm ref}\gets\Proj_{x_i}(r_r)/(m+\|\Proj_{x_i}(r_r)\|_{x_i})$, for $r\in\{0,1\}$
\State $v_i^{\rm ref}\gets b_{i0}q_{i0}^{\rm ref}+b_{i1}q_{i1}^{\rm ref}$
\State $v_{\theta,i}\gets v_i^{\rm part}+v_i^{\rm ref}\in T_{x_i}\HH_m$
\State \Return $V_\theta=(v_{\theta,1},\ldots,v_{\theta,N})$
\end{algorithmic}
\end{algorithm}

The signed logarithmic preprocessing in Algorithm~\ref{alg:network} is
\begin{equation}
 \psi(z)=\operatorname{sgn}(z)\log(1+|z|).
 \label{eq:signed-log}
\end{equation}
All learned maps $\phi_\bullet$ act on scalar invariant inputs. With
$\sigma=\operatorname{SiLU}$, their widths are
\begin{align}
 \phi_c &: d_c\to d_h\to d_h,& \phi_t &: 3\to d_h\to d_h,&
 \phi_{\rm seed}&:2\to d_h\to d_h,\\
 \phi_e^\ell &: (2d_h+2)\to d_h\to d_h,&
 \phi_h^\ell &:2d_h\to2d_h\to d_h,\\
 \phi_x^\ell &:d_h\to d_h\to1,&
 \phi_f &: (2d_h+2)\to d_h\to1,&
 \phi_r &: (d_h+2)\to d_h\to2.
 \label{eq:mlp-widths}
\end{align}
Every hidden linear layer is followed by $\sigma$; $\phi_e^\ell$ also applies
$\sigma$ after its second linear layer. The final layers of $\phi_x^\ell$, $\phi_f$,
and $\phi_r$ have weights initialized from $\mathcal{N}(0,10^{-6})$ and zero bias.
The auxiliary-vector residual scale is $\varepsilon_y=10^{-2}$.

There is no message gate, softmax, absolute-value operation, or positivity constraint.
Thus the graph aggregation and the final particle readout are signed sums. The
particle term is intrinsically tangent because $q_{ij}\in T_{x_i}\HH_m$, and the
reference term is tangent by Eq.~\eqref{eq:projection}. Therefore their sum is
the model output directly in $T_{x_i}\HH_m$.

For $\Lambda\in SO^+(1,3)$, the network is jointly Lorentz equivariant,
\begin{equation}
 V_\theta(\Lambda^{\times N}X,t,c;\Lambda r_0,\Lambda r_1)
 =\Lambda^{\times N}V_\theta(X,t,c;r_0,r_1).
 \label{eq:equivariance}
\end{equation}
The statement follows because MLP inputs are Lorentz scalars, learned vector
coefficients are scalars, and tangent maps commute with Lorentz transformations.
Holding the references fixed in the laboratory frame supplies conditional symmetry
breaking through the beam/time and jet directions.

\section{Mass-shell Riemannian flow matching}

\subsection{Source and data distributions}

Observed massless constituent momenta $\bar x_i=(\|\bm p_i\|_2,\bm p_i)$ are
lifted to the shell by preserving spatial momentum,
\begin{equation}
 \iota_m(\bar x_i)=\left(\sqrt{\|\bm p_i\|_2^2+m^2},\bm p_i\right).
 \label{eq:shell-lift}
\end{equation}
Let $X_1\sim q_1(\cdot\mid c)$ denote a lifted target jet. The conditional source
distribution $q_0(\cdot\mid c)$ is an axis-aligned lognormal prior. It samples
\begin{align}
 \Delta\eta_i,\Delta\phi_i&\stackrel{\rm iid}{\sim}\mathcal{N}(0,0.15^2),&
 z_i&\stackrel{\rm iid}{\sim}\mathcal{N}(0,1),\\
 w_i&=\frac{e^{z_i}}{\sum_j e^{z_j}},&
 \eta_i&=\eta_J+\Delta\eta_i,\qquad\phi_i=\phi_J+\Delta\phi_i.
\end{align}
Writing
\begin{equation}
 R=\left\|\sum_iw_i(\cos\phi_i,\sin\phi_i)\right\|_2,
 \qquad p_{T,i}=w_i\frac{p_{T,J}}{R},
 \label{eq:prior-pt}
\end{equation}
the unlifted source four-vector is
\begin{equation}
 \bar x_{0,i}=\frac{1}{s}
 (p_{T,i}\cosh\eta_i,\ p_{T,i}\cos\phi_i,\ p_{T,i}\sin\phi_i,\ p_{T,i}\sinh\eta_i),
\end{equation}
followed by the lift in Eq.~\eqref{eq:shell-lift}.

\subsection{Online geodesic coupling}

Particle indices do not identify physical objects. For every fresh source draw,
prior constituents are paired with target constituents through the per-jet
Hungarian assignment
\begin{equation}
 \pi^\star=\arg\min_{\pi\in S_N}\sum_{i=1}^N\dist(x_{0,\pi(i)},x_{1,i})^2.
 \label{eq:icp}
\end{equation}
The source is permuted according to $\pi^\star$ before the RFM path is formed.
This coupling is recomputed online for each training batch and does not reuse a
fixed noise-to-data pairing.

\subsection{Conditional geodesic path and objective}

For a coupled pair $(X_0,X_1)$ and $t\sim\mathcal{U}[0,1]$, RFM uses the
product-manifold geodesic
\begin{equation}
 X_t=\Exp_{X_1}\!\left((1-t)\Log_{X_1}(X_0)\right)
 =\Exp_{X_0}\!\left(t\Log_{X_0}(X_1)\right).
 \label{eq:rfm-path}
\end{equation}
Its conditional velocity, tangent at $X_t$, is
\begin{equation}
U_t(X_t\mid X_1)=\frac{\Log_{X_t}(X_1)}{1-t}\in T_{X_t}\mathcal{M}_N.
\label{eq:target-field}
\end{equation}
Equation~\eqref{eq:target-field} is the geodesic specialization of the
distance-gradient target in Chen and Lipman~\cite[Eqs.~(13)--(15)]{chen2023flow},
not a separate choice of target field.  Let $d_{\mathcal{M}}$ denote the geodesic
distance of the product metric.  Along the conditional geodesic,
\begin{equation}
 d_{\mathcal{M}}(X_t,X_1)=(1-t)d_{\mathcal{M}}(X_0,X_1),\qquad
 \|\nabla_{X_t}d_{\mathcal{M}}(X_t,X_1)\|_{X_t}=1,
 \label{eq:geodesic-distance-identities}
\end{equation}
and the standard distance-gradient identity is
\begin{equation}
 \Log_{X_t}(X_1)=-d_{\mathcal{M}}(X_t,X_1)
 \nabla_{X_t}d_{\mathcal{M}}(X_t,X_1).
 \label{eq:log-distance-gradient}
\end{equation}
For $\kappa(t)=1-t$, the Chen--Lipman target therefore reduces as
\begin{align}
 U_t(X_t\mid X_1)
 &=-d_{\mathcal{M}}(X_0,X_1)\nabla_{X_t}d_{\mathcal{M}}(X_t,X_1) \notag\\
 &=\frac{\Log_{X_t}(X_1)}{1-t}.
 \label{eq:target-field-derivation}
\end{align}
The numerator vanishes linearly as $t\to1$; implementation terminates just before
$t=1$ and uses a protected denominator.

The model is trained by the product-manifold squared error
\begin{equation}
 \mathcal{L}_{\rm RFM}(\theta)=
 \mathbb{E}\left[
 \left\|V_\theta(X_t,t,c;r_0,r_1)-U_t(X_t\mid X_1)\right\|_{X_t}^2
 \right],
 \label{eq:rfm-loss}
\end{equation}
where the norm is defined by Eq.~\eqref{eq:product-metric} and is averaged over
real particles in implementation. Under squared-error regression, the population
optimum is the marginal velocity
\begin{equation}
 V^\star(x,t,c)=\mathbb{E}[U_t(X_t\mid X_1)\mid X_t=x,c].
 \label{eq:marginal-field}
\end{equation}
This field has the same time marginals as the mixture of conditional geodesic paths,
so integrating it transports $q_0(\cdot\mid c)$ toward $q_1(\cdot\mid c)$.

\subsection{Sampling}

Generation starts from $X(0)\sim q_0(\cdot\mid c)$ and integrates
\begin{equation}
 \dot X(t)=V_\theta(X(t),t,c;r_0,r_1).
 \label{eq:ode}
\end{equation}
Using geodesic Euler integration, one step is
\begin{equation}
 X_{k+1}=\Exp_{X_k}\!\left(\Delta t\,V_\theta(X_k,t_k,c;r_0,r_1)\right).
 \label{eq:euler}
\end{equation}
The product exponential is applied constituentwise. Hence all generated particles
remain on the positive-energy mass shell throughout sampling. The reported Arm H
configuration uses $m=0.1$, 64 Euler steps, and an endpoint time of $0.99999$.

\section{Ablation results}
\label{sec:ablations}

Unless stated otherwise, the endpoint metrics below use EMA weights, 50,000
generated jets, and 64 geodesic Euler steps. FPND is computed after the validated
JetNet coordinate conversion; historical FPND values produced before that repair are
excluded. W1M, W1P, and FPD are reported in the native normalized JetNet convention.
These are single-run estimates, so small differences should not be interpreted as
seed-level significance.

\subsection{Geometry and readout construction}

The earliest controlled comparison replaced ambient Euclidean interpolation with the
regulated mass-shell path while holding the presence or absence of reference readout
fixed. Table~\ref{tab:geometry-ablation} reports only quantities that remain valid for
those historical runs. The mass-shell formulation reduced W1M by approximately an
order of magnitude and removed the large spacelike component population. Physical
support is therefore enforced by construction rather than learned as an approximate
constraint.

\begin{table}[H]
\centering
\caption{Early geometry ablation on gluon-30 at 200k updates. FPND is omitted because
these runs predate the corrected ParticleNet input convention. Percentages refer to
generated jets or generated constituents as appropriate.}
\label{tab:geometry-ablation}
\small
\begin{tabular}{clrrrr}
\toprule
Arm & Geometry and readout & W1M $\downarrow$ & Invalid (\%) & Spacelike (\%) & $E<0$ (\%) \\
\midrule
A & Euclidean, no references       & 0.249857 & 4.492 & 25.45 & 2.290 \\
B & Mass shell, no references      & 0.025672 & 0.000 &  0.00 & 0.000 \\
C & Euclidean, plain references    & 0.105699 & 0.006 & 33.87 & 0.298 \\
D & Mass shell, plain references   & 0.008915 & 0.004 &  0.00 & 0.000 \\
\bottomrule
\end{tabular}
\end{table}

The later G--J campaign isolates normalized tangent references versus raw reference
projections and evolving auxiliary geometry versus fixed physical geometry.
Table~\ref{tab:gj-ablation} shows a consistent advantage for evolving $y^\ell$ over
fixed geometry (G versus I and H versus J). Reference normalization is mildly favorable
in the evolving model but not decisive at this sample size. Reinjecting condition and
time embeddings in every block does not help.

\begin{table}[H]
\centering
\caption{Controlled gluon-30 construction ablations. ``Evol.'' denotes evolving
auxiliary Lorentz vectors and ``fixed'' denotes message geometry frozen to $x$.
The 500k rows test training duration rather than a new architecture.}
\label{tab:gj-ablation}
\resizebox{\linewidth}{!}{%
\begin{tabular}{llrrrrrr}
\toprule
Run & Construction & FPND$_g$ $\downarrow$ & FPD $\times10^4$ $\downarrow$ & W1M $\times10^3$ $\downarrow$ & mean W1P $\times10^3$ $\downarrow$ & Cov. $\uparrow$ & MMD $\times10^2$ $\downarrow$ \\
\midrule
G, 200k & raw references, evol.        & 0.3935 & 1.593 & 1.884 & 1.371 & 0.563 & 2.647 \\
H, 200k & normalized references, evol. & \textbf{0.3745} & \textbf{1.469} & \textbf{1.808} & \textbf{1.264} & \textbf{0.565} & 2.684 \\
I, 200k & raw references, fixed        & 0.4264 & 1.685 & 1.981 & 1.405 & 0.545 & 2.793 \\
J, 200k & normalized references, fixed & 0.4563 & 1.747 & 2.191 & 1.465 & 0.539 & 2.815 \\
H, 200k & per-block condition/time     & 0.3868 & 1.524 & 2.164 & 1.299 & 0.544 & 2.740 \\
G, 500k & raw references, evol.        & 0.3754 & \textbf{1.389} & \textbf{1.868} & 1.384 & \textbf{0.555} & \textbf{2.676} \\
H, 500k & normalized references, evol. & \textbf{0.3678} & 1.397 & 2.030 & \textbf{1.334} & 0.551 & 2.844 \\
\bottomrule
\end{tabular}}
\end{table}

\subsection{Coupling and the auxiliary workspace}

Table~\ref{tab:icp-ablation} is the matched 994k-update gluon-30 training comparison.
Removing online geodesic ICP increases FPND$_g$ from $0.391$ to $1.054$ and worsens
FPD, W1M, and W1P while retaining exact physical support. ICP therefore improves the
learned distribution rather than merely preventing invalid trajectories. Replacing the
physical log-map directions by raw latent displacements $y_j^L-y_i^L$ leaves the
already-degraded no-ICP result nearly unchanged. Because the completed latent run also
removes ICP, it does \emph{not} isolate the readout choice under the canonical coupling;
the corresponding ICP-enabled retraining remains necessary before simplifying the final
physical readout.

\begin{table}[H]
\centering
\caption{Matched 994k-update gluon-30 coupling/readout runs. ``Physical'' uses
coefficients from $(h^L,y^L)$ but directions from $\Log_{x_i}(x_j)$; ``latent'' uses
$y_j^L-y_i^L$ for the directions.}
\label{tab:icp-ablation}
\resizebox{\linewidth}{!}{%
\begin{tabular}{llrrrrrr}
\toprule
Coupling & Directions & FPND$_g$ $\downarrow$ & FPD $\times10^4$ $\downarrow$ & W1M $\times10^3$ $\downarrow$ & mean W1P $\times10^3$ $\downarrow$ & Cov. $\uparrow$ & MMD $\times10^2$ $\downarrow$ \\
\midrule
online ICP & physical & \textbf{0.3910} & \textbf{1.448} & \textbf{2.025} & \textbf{1.529} & \textbf{0.556} & 2.690 \\
none       & physical & 1.0541 & 4.859 & 3.932 & 2.763 & 0.536 & \textbf{2.630} \\
none       & latent   & 1.0459 & 4.837 & 3.918 & 2.825 & 0.553 & 2.758 \\
\bottomrule
\end{tabular}}
\end{table}

\subsection{Rejected auxiliary mechanisms and capacity reduction}

The matched 100k-update Gram-auxiliary experiment in
Table~\ref{tab:gram-ablation} improves a few one-dimensional quantities but damages the
joint 36-EFP FPD by a factor of $4.3$ and worsens W1P. It was therefore rejected. The
2k token smoke tests in Table~\ref{tab:token-ablation} likewise show no loss advantage;
the scalar-vector token violates the endpoint-invalidity gate. Neither mechanism is
present in the final model.

\begin{table}[H]
\centering
\caption{Matched gluon-30 auxiliary-objective ablation at 100k updates. Historical
FPND is omitted because these runs predate the corrected input convention.}
\label{tab:gram-ablation}
\resizebox{\linewidth}{!}{%
\begin{tabular}{lrrrrrr}
\toprule
Objective & W1M $\times10^3$ & mean W1P $\times10^3$ & mean W1EFP $\times10^5$ & FPD $\times10^4$ & Cov. & MMD $\times10^2$ \\
\midrule
control & 1.805 & \textbf{0.959} & 1.487 & \textbf{6.15} & 0.54 & 2.681 \\
Gram-50 & \textbf{1.724} & 1.137 & \textbf{1.323} & 26.40 & \textbf{0.56} & \textbf{2.619} \\
\bottomrule
\end{tabular}}
\end{table}

\begin{table}[H]
\centering
\caption{Jet-token qualification smoke test at 2k optimizer updates.}
\label{tab:token-ablation}
\small
\begin{tabular}{lrrl}
\toprule
Token & Final median loss $\downarrow$ & Invalid endpoints & Qualification \\
\midrule
none          & \textbf{0.004798} & 0/2000 & pass \\
scalar        & 0.005057 & 1/2000 & pass with warning \\
scalar+vector & ---      & 3/2000 & fail \\
\bottomrule
\end{tabular}
\end{table}

The subsequent backbone simplification reduced parameter count by more than an order of
magnitude without a corresponding collapse in corrected FPND. Table~\ref{tab:capacity}
is a capacity frontier rather than a perfectly matched width ablation: variants a--c use
width 96, d uses width 128, and the historical attention model uses width 256. The
readout-only, no-pooling variant is the smallest and best of the message-passing variants
on FPND; global pooling is not justified by these results.

\begin{table}[H]
\centering
\caption{Gluon-30 capacity and message-passing simplification at 600k updates. The
attention-model FPND was recomputed after the coordinate repair; its other historical
metrics are omitted here.}
\label{tab:capacity}
\resizebox{\linewidth}{!}{%
\begin{tabular}{lrrrrrr}
\toprule
Backbone & Parameters & FPND$_g$ $\downarrow$ & FPD $\times10^4$ $\downarrow$ & W1M $\times10^3$ $\downarrow$ & Cov. $\uparrow$ & MMD $\times10^2$ $\downarrow$ \\
\midrule
tangent attention, width 256       & 7.39M & \textbf{0.507} & --- & --- & --- & --- \\
readout-only, no pool, width 96    & 0.513M & 0.5295 & 3.484 & 2.431 & 0.550 & \textbf{2.591} \\
tangent channels, no pool, width 96& 0.560M & 0.5504 & 2.995 & \textbf{2.310} & 0.536 & 2.744 \\
readout-only, global pool, width 96 & 0.848M & 0.5460 & 2.760 & 2.602 & \textbf{0.552} & 2.704 \\
tangent channels, global pool, width 128 & 0.910M & 0.553 & \textbf{2.462} & 2.635 & 0.548 & 2.704 \\
\bottomrule
\end{tabular}}
\end{table}

\subsection{Training stability and readout initialization}

Zero-initializing the terminal velocity heads makes the initial sampler the identity map.
The 2k-update qualification in Table~\ref{tab:init-ablation} shows that initialization
strongly changes the route to a stable solution, but the terminal losses and endpoint
validity are tied. Zero initialization was retained as the safer default; this is not a
claim of better converged sample quality. The experiment also demonstrates why loss
alone is insufficient for checkpoint selection: either arm can have an ordinary loss
while fixed-seed sampler probes remain unstable.

\begin{table}[H]
\centering
\caption{Velocity-readout initialization qualification. Probe entries report unstable
trajectories out of 64.}
\label{tab:init-ablation}
\small
\begin{tabular}{lrr}
\toprule
Quantity & small-normal & zero \\
\midrule
unstable at step 0      & 64/64 & \textbf{0/64} \\
unstable at step 100    & 64/64 & \textbf{0/64} \\
unstable at step 500    & \textbf{60/64} & 63/64 \\
unstable at step 1000   & \textbf{0/64} & 18/64 \\
unstable at step 2000   & 0/64 & 0/64 \\
first-200 median loss   & 7.664803 & \textbf{0.015781} \\
final-200 median loss   & 0.004732 & \textbf{0.004675} \\
final failures (of 128) & 0 & 0 \\
\bottomrule
\end{tabular}
\end{table}

\subsection{Inference-time simplifications and scaling}

The GQT-994k checkpoint gives a clean integration-resolution curve without retraining.
As shown in Table~\ref{tab:euler-ablation}, doubling from 32 to 64 and from 64 to 128
steps improves FPND in every class while approximately doubling generation time. The
continued gain at 128 steps indicates discretization error at the canonical 64-step
setting rather than a saturated sampler.

\begin{table}[H]
\centering
\caption{Geodesic Euler scaling for the GQT-994k checkpoint at fixed $w=0$. Generation
time is for 150,000 class-stratified samples on the same A40 inference setup.}
\label{tab:euler-ablation}
\small
\begin{tabular}{rrrrr}
\toprule
Steps & FPND$_g$ $\downarrow$ & FPND$_q$ $\downarrow$ & FPND$_t$ $\downarrow$ & Generation time (s) \\
\midrule
32  & 0.4059 & 0.6132 & 0.9646 &  992.1 \\
64  & 0.1420 & 0.2679 & 0.5244 & 1853.8 \\
128 & \textbf{0.0797} & \textbf{0.1546} & \textbf{0.3693} & 3713.5 \\
\bottomrule
\end{tabular}
\end{table}

The class-conditional guidance sweep in Table~\ref{tab:cfg-sweep} is sharply
class-dependent. Guidance is unnecessary for gluons, while light-quark and top jets
prefer $w=0.25$ and $w=0.5$, respectively. Larger weights rapidly degrade all classes.

\begin{table}[H]
\centering
\caption{Classwise FPND for the GQT-994k classifier-free-guidance sweep at 64 Euler
steps. Bold entries are the selected per-class deployment settings.}
\label{tab:cfg-sweep}
\small
\begin{tabular}{rrrr}
\toprule
$w$ & FPND$_g$ $\downarrow$ & FPND$_q$ $\downarrow$ & FPND$_t$ $\downarrow$ \\
\midrule
0.00 & \textbf{0.1420} & 0.2679 & 0.5244 \\
0.25 & 0.3182 & \textbf{0.2137} & 0.2801 \\
0.50 & 0.5255 & 0.3138 & \textbf{0.2019} \\
0.75 & 0.7063 & 0.5796 & 0.2582 \\
1.00 & 0.9160 & 1.6340 & 0.4261 \\
2.00 & 4.5388 & 5.7217 & 2.1008 \\
3.00 & 25.0270 & 15.1774 & 8.4812 \\
\bottomrule
\end{tabular}
\end{table}

Finally, Table~\ref{tab:projection-ablation} tests the redundant terminal tangent
projection that previously followed the sum of already-tangent particle and reference
terms. Turning it off leaves FPND unchanged to approximately six significant figures,
and the 36-EFP FPD is numerically identical. The projection was therefore deleted from
both implementation and algorithm.

\begin{table}[H]
\centering
\caption{Final-projection ablation on the GQT-994k checkpoint at the selected per-class
guidance weights. ``On'' applies $\Pi_{x_i}$ to the final sum; ``off'' returns that
already-tangent sum directly.}
\label{tab:projection-ablation}
\small
\begin{tabular}{lrrrr}
\toprule
Class ($w$) & FPND on & FPND off & FPD on $\times10^4$ & FPD off $\times10^4$ \\
\midrule
g ($0$)    & 0.14199731 & 0.14199728 & 8.7697 & 8.7697 \\
q ($0.25$) & 0.21368916 & 0.21369031 & 33.2823 & 33.2823 \\
t ($0.5$)  & 0.20192473 & 0.20192409 & 24.2870 & 24.2870 \\
\bottomrule
\end{tabular}
\end{table}

\begin{thebibliography}{9}
\bibitem{lou2020neural}
Aaron Lou, Derek Lim, Isay Katsman, Leo Huang, Qingxuan Jiang, Ser-Nam Lim, and
Christopher M. De Sa.
\newblock Neural Manifold Ordinary Differential Equations.
\newblock In \emph{Advances in Neural Information Processing Systems}, 2020.
\newblock Appendix B.4, Table 2.

\bibitem{chen2023flow}
Ricky T. Q. Chen and Yaron Lipman.
\newblock Flow Matching on General Geometries.
\newblock \emph{arXiv preprint arXiv:2302.03660}, 2023.
\end{thebibliography}

\end{document}
